\subsection{子集集合公理}

\[
(\forall \mathbf {X}) \operatorname{Coll} _ {\mathbf {Y}} (\mathbf {Y} \subset \mathbf {X}).
\]

该公理意味着，对于每个集合 X，存在一个集合，其元素是 X 的所有子集，即集合 \(\mathcal{E}_{\mathbf{Y}}(\mathbf{Y}\subset\mathbf{X})\) （§1，第4节）；该集合记为 \(\mathfrak{P}(\mathbf{X})\) 并被称为X的子集集合。显然，如果 \(X\subset X'\) ，我们有 \(\mathfrak{P}(\mathbf{X})\subset\mathfrak{P}(\mathbf{X}')\) .

设 A、B 为两个集合， \(\Gamma\) 为 A 与 B 之间的对应。函数 \(X \to \Gamma\langle X\rangle\) ( \(X \subset \mathfrak{P}(A)\) , \(\Gamma\langle X\rangle \in \mathfrak{P}(B)\) ) 称为 \(\Gamma\) 到子集集合上的典范扩张（或简称扩张），记作 \(\hat{\Gamma}\) ；它是从 \(\mathfrak{P}(A)\) 到 \(\mathfrak{P}(B)\) 的映射。如果 \(\Gamma'\) 是 B 与集合 C 之间的一个对应，则公式 \((\Gamma' \circ \Gamma)\langle X\rangle = \Gamma'\langle\Gamma\langle X\rangle\rangle\) 表明， \(\Gamma' \circ \Gamma\) 到子集集合上的扩张就是映射 \(\hat{\Gamma}' \circ \hat{\Gamma}\) .

命题1. (1) 若 \(f\) 是集合 \(\mathbf{E}\) 到集合 \(\mathbf{F}\) 的满射，则典范扩张 \(\hat{f}\) （ \(f\) 的）是从 \(\mathfrak{P}(\mathbf{E})\) 到 \(\mathfrak{P}(\mathbf{F})\) 的满射。

\begin{enumerate}
\def\labelenumi{(\arabic{enumi})}
\setcounter{enumi}{1}
\item
  如果 \(f\) 是从 \(\mathbf{E}\) 到 \(\mathbf{F}\) 的单射，则典范扩张 \(\hat{f}\) （ \(f\) 的）是从 \(\mathfrak{P}(\mathbf{E})\) 到 \(\mathfrak{P}(\mathbf{F})\) 的单射。
\item
  如果 \(s\) 是 \(f\) ，那么 \(f \circ s\) 是 \(\mathbf{F}\) ，因此 \(\hat{f} \circ \hat{s}\) 是 \(\mathfrak{P}(\mathbf{F})\) 中；因此 \(\hat{f}\) 是满射，且 \(\hat{s}\) 是 \(\hat{f}\) （§3，第8节）。
\item
  该命题在以下情形下是显然的： \(\mathbf{E} = \varnothing\) ，因为那样 \(\mathfrak{P}(\mathbf{E}) = \{\varnothing\}\) 。如果 \(\mathbf{E} \neq \varnothing\) 并且如果 \(r\) 是 \(f\) ，那么 \(r \circ f\) 是 \(\mathbf{E}\) ，因此 \(\hat{r} \circ \hat{f}\) 是 \(\mathfrak{P}(\mathbf{E})\) 中；因此 \(\hat{f}\) 是单射，并且 \(\hat{r}\) 是 \(\hat{f}\) （§3，第8节）。
\end{enumerate}

